% =============================================================================
% Relatorio de andamento: replicacao de An et al. (2018) e analise de estabilidade
% Figuras e tabelas sao geradas por:  python scripts/report_assets.py [...]
% (ver secao "Como conferir os resultados"). Nada aqui e digitado a partir de
% graficos: os numeros vem dos arquivos de resultado.
% =============================================================================
\documentclass[11pt,a4paper]{article}

\usepackage[utf8]{inputenc}
\usepackage[T1]{fontenc}
\usepackage[brazil]{babel}
\usepackage{lmodern}
\usepackage[a4paper,margin=2.5cm]{geometry}
\usepackage{amsmath,amssymb}
\usepackage{graphicx}
\usepackage{booktabs}
\usepackage{array}
\usepackage{xcolor}
\usepackage[font=small,labelfont=bf]{caption}
\usepackage{enumitem}
\usepackage[section]{placeins}   % figuras nao saem da propria secao
\usepackage{xurl}
\usepackage[hidelinks]{hyperref}
\usepackage{microtype}

\setlist{itemsep=2pt,topsep=4pt}
\setlength{\parskip}{4pt}

\definecolor{pend}{HTML}{8A8985}
\definecolor{pendfundo}{HTML}{F3F2EF}

% Caixa cinza para o que ainda esta rodando
\newcommand{\pendente}[1]{%
  \par\noindent\colorbox{pendfundo}{\parbox{\dimexpr\linewidth-2\fboxsep}{%
    \small\color{pend}\textbf{Pendente:} #1}}\par}

% Inclui a figura se ela ja foi gerada; senao mostra o que falta rodar
\newcommand{\figura}[3]{%
  \IfFileExists{figuras/#1.pdf}%
    {\includegraphics[width=#2\linewidth]{figuras/#1.pdf}}%
    {\pendente{figura \texttt{\detokenize{#1}} ainda não gerada. #3}}}

\title{\textbf{Leis de aprendizado sob a ótica de controle:\\
replicação de An et al.\ (2018) e análise de estabilidade em tempo discreto}\\[4pt]
\large Relatório de andamento --- Projeto Final de Curso I}
\author{Laura Inês Freitas Campos\\
\small Orientador: Prof.\ Dr.\ Everthon de Souza Oliveira\\
\small Departamento de Engenharia Elétrica --- CEFET-MG}
\date{\today}

\begin{document}
\maketitle

\begin{abstract}
\noindent
Este relatório descreve a replicação dos resultados de An et al.\ (2018), que
propõem tratar o treinamento de redes neurais como uma malha de controle e
acrescentar um termo derivativo (D) ao gradiente descendente com momento.
Replicamos o experimento do artigo com a rede MLP no MNIST, com 10 repetições, e o
exemplo de duas variáveis usado pelos autores para ilustrar o efeito do termo D.
A vantagem do método sobre o momento foi reproduzida, mas ela praticamente não
depende do termo D. Encontramos também diferenças entre o texto do artigo e o
código publicado pelos autores, entre elas o \emph{sinal} do termo D. Para explicar
esses resultados, escrevemos cada lei de aprendizado como um sistema linear em
tempo discreto. A posição dos polos previu corretamente, nos 14 casos testados,
quando o treinamento converge e quando diverge.
\end{abstract}

% =============================================================================
\section{Objetivo}

O orientador pediu a replicação de An et al.\ (2018) como ponto de partida do
PFC. O objetivo foi verificar três coisas:
\begin{enumerate}
  \item se os resultados do artigo se reproduzem com o código e os dados disponíveis;
  \item se o texto do artigo e o código público descrevem o mesmo algoritmo;
  \item se a leitura ``otimizador como controlador'' explica o comportamento
        observado, que é o tema central do PFC (objetivo específico b da proposta).
\end{enumerate}

% =============================================================================
\section{A ideia, em linguagem simples}

Treinar uma rede neural é ajustar seus pesos para diminuir o erro. A cada passo,
o \textbf{gradiente} indica para que lado o erro aumenta, e o algoritmo dá um passo
no sentido contrário. O tamanho desse passo é a \textbf{taxa de aprendizado}.

Visto como engenheiro de controle, isso é uma \textbf{malha fechada}: o gradiente
faz o papel do \emph{sinal de erro}, a regra que transforma o gradiente em passo
faz o papel do \emph{controlador}, e a função de erro da rede é a \emph{planta}. Daí
a correspondência proposta por An et al.:

\begin{center}
\begin{tabular}{llp{6.2cm}}
\toprule
Algoritmo & Ação de controle & O que o passo usa \\
\midrule
SGD (gradiente descendente) & P (proporcional) & só o gradiente atual \\
SGD com momento & PI (prop.-integral) & também a soma ponderada dos gradientes passados \\
PID de An et al. & PID & também a \emph{variação} do gradiente \\
\bottomrule
\end{tabular}
\end{center}

O argumento do artigo: o momento, por acumular o passado (ação integral), passa do
ponto e oscila em torno do mínimo, assim como um controlador PI com ganho
integral alto. Um termo derivativo funcionaria como um \textbf{amortecedor},
reduzindo essa oscilação, que os autores chamam de \emph{overshoot} (sobressinal).

% =============================================================================
\section{O que foi feito}

\begin{itemize}
  \item \textbf{Ambiente de experimentos.} Repositório público
        (\url{https://github.com/laurafreitascampos/otimizadores-controle}) em
        Python/PyTorch, que aproveita os otimizadores do trabalho do Laboratório de IA.
        Cada experimento é descrito por um arquivo de configuração. Cada
        execução grava num arquivo \texttt{results.json} a configuração usada, a
        versão exata do código (\emph{commit} do git), o hardware e as curvas de
        cada repetição.
  \item \textbf{Testes de equivalência.} Antes de qualquer experimento, verificamos
        automaticamente que o PID com só o termo P reproduz o SGD do PyTorch, e
        que com só o termo I reproduz o SGD com momento, passo a passo, com
        diferença menor que $10^{-10}$. Esses testes confirmam no código as
        afirmações das Seções~3.2 e 3.3 do artigo.
  \item \textbf{Protocolo justo.} Todos os algoritmos comparados partem dos
        \emph{mesmos pesos iniciais} e veem os dados \emph{na mesma ordem}. Cada
        comparação é repetida com 10 \textbf{sementes} (cada semente fixa um sorteio
        diferente de pesos iniciais e de ordem dos dados), e reportamos
        média~$\pm$~desvio padrão.
  \item \textbf{Reprodutibilidade.} O experimento de duas variáveis deu resultados
        idênticos, até a última casa decimal, em Windows (com placa de vídeo) e em Linux.
\end{itemize}

% =============================================================================
\section{Resultado 1: o exemplo de duas variáveis e o sinal do termo D}

Os autores ilustram o efeito do termo D (Figura~2 do artigo) numa ``tigela''
alongada, $f(x_1,x_2) = 0{,}1\,x_1^2 + 2\,x_2^2$, partindo de $(-10,-10)$. A direção
$x_2$ é íngreme (curvatura alta) e é onde o momento oscila. Rodamos o mesmo
exemplo com o termo D escrito \emph{exatamente como na Eq.~(10) do artigo} e
também com o sinal trocado (Fig.~\ref{fig:toy}).

\begin{figure}[!htbp]
  \centering
  \figura{toy_resposta_x2}{1.0}{Rode \texttt{python scripts/report\_assets.py}.}
  \caption{Resposta no tempo da coordenada $x_2$ no exemplo de duas variáveis.
  (a) Com o termo D somado, como escrito na Eq.~(10), o sobressinal \emph{aumenta}
  e, para $K_d \geq 0{,}5$, o algoritmo diverge; nesses casos,
  já no primeiro passo ele se afasta do mínimo. (b) Com o sinal trocado, o termo D amortece a oscilação:
  com $K_d = -0{,}93$ o sobressinal desaparece. Taxa de aprendizado
  $0{,}01$ e momento $0{,}9$; o artigo não informa esses valores para esse exemplo.}
  \label{fig:toy}
\end{figure}

\textbf{O que isso mostra.} O comportamento da Figura~2 do artigo (menos
sobressinal com o termo D) só aparece com o sinal \emph{oposto} ao da Eq.~(10).
É exatamente o sinal que o código público dos autores usa: lá o termo D entra
\emph{subtraindo}, junto com o gradiente. Em linguagem de controle, o amortecedor
precisa empurrar \emph{contra} o movimento; com o sinal do texto, ele empurra
\emph{a favor} e injeta energia no sistema.

% =============================================================================
\section{Resultado 2: replicação da Figura 4 (MLP no MNIST)}

\textbf{Configuração} (igual à do artigo): rede com uma camada oculta de 1000
neurônios, base MNIST de dígitos manuscritos, 20 épocas (uma \textbf{época} é uma
passada completa pelos dados de treino), taxa de aprendizado $0{,}01$, momento
$0{,}9$. Comparamos três algoritmos:
\begin{itemize}
  \item o \textbf{SGD com momento} dos autores;
  \item o \textbf{PID como no código dos autores} ($K_i = 3$, $K_d = 100$);
  \item o \textbf{mesmo PID com $K_d = 0$}, isto é, sem o termo D. Este último é um
        controle que acrescentamos para isolar o efeito do termo D.
\end{itemize}

\IfFileExists{tabelas/fig4.tex}{%
\begin{table}[!htbp]
  \centering\small
  \caption{Replicação da Figura~4: acurácia no conjunto de teste após 20 épocas e
  número de épocas até atingir 98\,\%. Média $\pm$ desvio padrão amostral entre
  sementes; $\Delta$ e intervalo de confiança de 95\,\% calculados semente a semente
  (pareado), contra o SGD com momento.}
  \label{tab:fig4}
  \input{tabelas/fig4.tex}
\end{table}
}{%
\begin{table}[!htbp]
  \centering\small
  \caption{Replicação da Figura~4 (10 sementes, 20 épocas, avaliação no conjunto de
  teste). Valores transcritos do resumo impresso pelo \texttt{train.py};
  desvio padrão com divisor $n$. A tabela com testes estatísticos será gerada
  a partir dos arquivos de resultado.}
  \label{tab:fig4}
  \begin{tabular}{llcc}
    \toprule
    Lote & Algoritmo & Acurácia final (\%) & Épocas até 98\,\% \\
    \midrule
    100 (código) & SGD com momento & $98{,}25 \pm 0{,}07$ & $6{,}8$ \\
                 & PID ($K_d = 100$) & $98{,}44 \pm 0{,}03$ & $4{,}1$ \\
                 & PID sem D ($K_d = 0$) & $98{,}43 \pm 0{,}06$ & $4{,}4$ \\
    \midrule
    128 (artigo) & SGD com momento & $98{,}23 \pm 0{,}06$ & $8{,}6$ \\
                 & PID ($K_d = 100$) & $98{,}40 \pm 0{,}05$ & $4{,}3$ \\
                 & PID com $K_d$ da Eq.~(14) & \multicolumn{2}{c}{divergiu nas 10 sementes} \\
    \midrule
    \multicolumn{2}{l}{\emph{Registro dos autores (1 execução, lote 100)}} & & \\
                 & SGD com momento & $97{,}99$ & --- \\
                 & PID & $98{,}26$ & --- \\
    \bottomrule
  \end{tabular}
\end{table}
}

\begin{figure}[!htbp]
  \centering
  \figura{fig4_mnist}{1.0}{Rode \texttt{report\_assets.py -{}-fig4 <results.json>}
  com o resultado da Fig.~4.}
  \caption{Replicação da Figura~4. (a) Acurácia de teste; (b) erro de treino, em
  escala logarítmica. Linha: média; faixa: $\pm$ um desvio padrão entre sementes.}
  \label{fig:fig4}
\end{figure}

\textbf{O que isso mostra.}
\begin{itemize}
  \item \textbf{A vantagem do PID se reproduz.} Ele termina com cerca de
        0,2 ponto percentual a mais de acurácia e chega a 98\,\% em cerca de 4
        épocas, contra 7 a 9 do momento. Nossos valores ficam ligeiramente acima
        dos registros dos autores, e com 10 repetições em vez de uma.
  \item \textbf{Mas a vantagem não vem do termo D.} Sem o termo D, o resultado
        é praticamente o mesmo (98,44\,\% contra 98,43\,\%). A diferença é muito
        menor que a variação entre sementes.
  \item \textbf{O ajuste de $K_d$ proposto no texto não funciona como escrito.} O
        valor dado pela Eq.~(14) para taxa de aprendizado $0{,}01$ é
        $K_d \approx 1855$, enquanto o código usa um $K_d$ efetivo de~1. Aplicado na
        forma da Eq.~(10), esse valor diverge em todas as sementes.
  \item \textbf{A métrica de ``épocas'' depende do tamanho do lote.} O mesmo
        momento leva 6,8 épocas com lote 100 e 8,6 com lote 128, porque lotes maiores
        significam menos passos por época. Comparações em épocas precisam fixar o lote.
\end{itemize}

% =============================================================================
\section{Resultado 3: análise de estabilidade em tempo discreto}

\subsection{A ideia}

Perto de um mínimo, o erro da rede se parece com uma tigela: em cada direção,
$f \approx \tfrac{1}{2}\,a\,\theta^2$, em que $a$ é a \textbf{curvatura} daquela
direção (quão íngreme ela é). O gradiente é então $a\,\theta$, e cada lei de
aprendizado vira um \textbf{sistema linear em tempo discreto}
\[
  \mathbf{x}_{t+1} = A\,\mathbf{x}_t ,
\]
em que o vetor $\mathbf{x}$ guarda o peso e as ``memórias'' do algoritmo (a
velocidade do momento, os termos I e D). Como em qualquer sistema discreto, quem
decide o comportamento são os \textbf{polos} de $A$ (seus autovalores):
\begin{itemize}
  \item todos dentro do círculo unitário ($|z|<1$): converge;
  \item algum fora ($|z|>1$): diverge;
  \item polos complexos: converge oscilando, com sobressinal.
\end{itemize}
Resumimos tudo no \textbf{raio espectral} $\rho$, o maior $|z|$: quanto menor,
mais rápido converge. Para o SGD com momento, por exemplo,
\[
  A = \begin{bmatrix} 1-h & \alpha \\ -h & \alpha \end{bmatrix},
  \qquad z^2 - (1+\alpha-h)\,z + \alpha = 0,
  \qquad h = \text{(taxa de aprendizado)} \times a .
\]
Essas matrizes são \emph{exatas} para funções quadráticas, e não uma aproximação. Um
teste automático confirma que elas reproduzem, passo a passo, os mesmos
algoritmos usados no MNIST.

\subsection{Os regimes do momento}

\begin{figure}[!htbp]
  \centering
  \figura{regimes_momentum}{1.0}{Rode \texttt{python scripts/report\_assets.py}.}
  \caption{Regimes do SGD com momento em função de $h$ (passo $\times$ curvatura) e do
  coeficiente de momento $\alpha$. Fronteiras calculadas analiticamente:
  $h = (1\mp\sqrt{\alpha})^2$ separa polos reais de complexos e $h = 2(1+\alpha)$ é o
  limite de estabilidade. No exemplo de duas variáveis, $x_1$ cai na região sem
  oscilação e $x_2$ na região oscilatória, como observado.}
  \label{fig:regimes}
\end{figure}

A Fig.~\ref{fig:regimes} explica por que, no exemplo de duas variáveis, só a
direção $x_2$ oscila: com $\alpha = 0{,}9$, a direção $x_1$ (curvatura $0{,}2$) fica
na região de polos reais, onde não há sobressinal, e a direção $x_2$ (curvatura 4)
fica na região de polos complexos.

\subsection{O PID do artigo: teoria contra experimento}

\begin{figure}[!htbp]
  \centering
  \figura{mapa_pid_artigo}{1.0}{Rode \texttt{python scripts/report\_assets.py}.}
  \caption{Estabilidade do PID da Eq.~(10) em função da curvatura da direção e de
  $K_d$ (taxa de aprendizado $0{,}01$, momento $0{,}9$). Cinza: região instável
  prevista ($\rho \geq 1$). Linha tracejada: onde o termo D não muda a velocidade de
  convergência; ela separa as regiões em que o termo D acelera ou atrasa em
  relação ao momento. Os marcadores são os 14 casos do exemplo de duas variáveis.}
  \label{fig:mapa}
\end{figure}

\begin{table}[!htbp]
  \centering\small
  \caption{Previsão pelo raio espectral $\rho$ contra o comportamento observado no
  exemplo de duas variáveis. Acomodação: número de passos até ficar a menos de 2\,\%
  da distância inicial ao mínimo.}
  \label{tab:validacao}
  \IfFileExists{tabelas/validacao_toy.tex}{% gerado por scripts/report_assets.py (determinístico)
\begin{tabular}{ccrcccrr}
\toprule
Coord. & $a$ & $K_d$ & $\rho$ previsto & Polos & Previsão & Sobressinal & Acomodação \\
\midrule
$x_1$ & $0{,}2$ & $+0{,}00$ & $0{,}9735$ & reais & converge \checkmark & $0{,}00$ & 161 \\
$x_1$ & $0{,}2$ & $+0{,}10$ & $0{,}9724$ & reais & converge \checkmark & $0{,}00$ & 157 \\
$x_1$ & $0{,}2$ & $+0{,}50$ & $0{,}9648$ & reais & converge \checkmark & $0{,}00$ & 139 \\
$x_1$ & $0{,}2$ & $+0{,}93$ & $0{,}9584$ & complexos & converge \checkmark & $0{,}00$ & 118 \\
$x_1$ & $0{,}2$ & $-0{,}10$ & $0{,}9745$ & reais & converge \checkmark & $0{,}00$ & 165 \\
$x_1$ & $0{,}2$ & $-0{,}50$ & $0{,}9777$ & reais & converge \checkmark & $0{,}00$ & 180 \\
$x_1$ & $0{,}2$ & $-0{,}93$ & $0{,}9801$ & reais & converge \checkmark & $0{,}00$ & 195 \\
$x_2$ & $4{,}0$ & $+0{,}00$ & $0{,}9487$ & complexos & converge \checkmark & $0{,}44$ & 68 \\
$x_2$ & $4{,}0$ & $+0{,}10$ & $0{,}9695$ & complexos & converge \checkmark & $0{,}62$ & 128 \\
$x_2$ & $4{,}0$ & $+0{,}50$ & $1{,}0488$ & complexos & diverge \checkmark & --- & --- \\
$x_2$ & $4{,}0$ & $+0{,}93$ & $1{,}1278$ & complexos & diverge \checkmark & --- & --- \\
$x_2$ & $4{,}0$ & $-0{,}10$ & $0{,}9274$ & complexos & converge \checkmark & $0{,}31$ & 51 \\
$x_2$ & $4{,}0$ & $-0{,}50$ & $0{,}9000$ & complexos & converge \checkmark & $0{,}07$ & 23 \\
$x_2$ & $4{,}0$ & $-0{,}93$ & $0{,}9038$ & reais & converge \checkmark & $0{,}00$ & 8 \\
\bottomrule
\end{tabular}
}%
    {\pendente{rode \texttt{python scripts/report\_assets.py}.}}
\end{table}

\textbf{O que isso mostra.} A teoria acerta se o algoritmo converge ou diverge em
\textbf{14 de 14} casos (Tab.~\ref{tab:validacao}), e explica dois fatos que só olhar
o gráfico não explica:
\begin{itemize}
  \item na direção de curvatura alta ($x_2$, polos complexos), o termo D com o sinal
        trocado amortece a oscilação. Com o sinal do artigo, ele empurra os polos
        para fora do círculo unitário ($\rho = 1{,}049$ e $1{,}128$), e o algoritmo
        diverge;
  \item na direção de curvatura baixa ($x_1$, polos reais), o mesmo termo D que
        ajuda em $x_2$ \emph{atrasa} a convergência: de 161 para 195 passos. Amortecer
        um sistema que já não oscila só o deixa mais lento.
\end{itemize}
\textbf{Consequência para o projeto:} não existe um único $K_d$ que seja bom para
todas as direções. Uma rede neural tem milhões de direções com curvaturas muito
diferentes; esse é um primeiro indicativo de projeto para o PFC.

\FloatBarrier
\subsection{Margem de estabilidade do PID usado no código}

\begin{figure}[!htbp]
  \centering
  \figura{margem_pid_codigo}{1.0}{Rode \texttt{python scripts/report\_assets.py}.}
  \caption{Maior valor de $h$ (passo $\times$ curvatura) para o qual o PID do código
  dos autores ainda é estável, em função de $K_d$. Para comparação, o SGD com momento
  ($\alpha = 0{,}9$) é estável até $h = 3{,}8$.}
  \label{fig:margem}
\end{figure}

A Fig.~\ref{fig:margem} mostra o preço do PID do código:
\begin{itemize}
  \item mesmo com $K_d = 0$, ele só é estável até $h \approx 0{,}78$, cerca de cinco
        vezes menos margem que o momento. O motivo é o ganho integral maior ($K_i = 3$);
  \item com $K_d = 100$, valor dos autores, a margem cai para $h \approx 0{,}15$;
  \item com $K_d$ negativo, o sinal do texto do artigo, a margem cai ainda mais rápido.
\end{itemize}
Em termos de projeto, o termo D \emph{consome} margem de estabilidade.

% =============================================================================
\section{Diferenças entre o artigo e o código público}

\begin{table}[!htbp]
  \centering\small
  \caption{Discrepâncias encontradas durante a replicação.}
  \label{tab:discrepancias}
  \begin{tabular}{>{\raggedright}p{2.6cm}>{\raggedright}p{4.3cm}>{\raggedright}p{3.6cm}p{3.6cm}}
    \toprule
    Item & Artigo & Código público & Observado aqui \\
    \midrule
    Forma da atualização & Eq.~(10): $\theta \leftarrow \theta + V + K_d D$ &
      $\theta \leftarrow \theta - \text{lr}\,(g + K_i I + K_d D)$ &
      equivalentes, a menos do sinal de $K_d$ (teste automático) \\
    Sinal do termo D & soma & subtrai & só o sinal do código amortece (Fig.~\ref{fig:toy}) \\
    Valor de $K_d$ & Eq.~(14): $\approx 1855$ para lr $=0{,}01$ & efetivo $= 1$ &
      Eq.~(14) diverge nas 10 sementes \\
    Tamanho do lote & 128 & 100 & rodamos os dois \\
    Decaimento de peso & não mencionado & $10^{-4}$ & usado nos dois casos \\
    Sobressinal, Eq.~(8) & dividido por $\theta^*$ & --- & indefinido com $\theta^* = 0$; normalizamos pela distância inicial \\
    Material suplementar & citado & --- & não disponível \\
    \bottomrule
  \end{tabular}
\end{table}

% =============================================================================
\section{Em andamento}

\begin{itemize}
  \item \textbf{De onde vem a vantagem do PID?} Com um gradiente constante, o PID
        do código ($K_i = 3$) dá passos 3,1 vezes maiores que o momento de mesma
        taxa de aprendizado: o ganho DC é $1 + 3/(1-0{,}9) = 31$, contra
        $1/(1-0{,}9) = 10$. Hipótese: o ganho no MNIST vem desse passo maior, e não da
        ação PID. Teste: o SGD com momento e taxa de aprendizado $0{,}031$, que tem o
        mesmo ganho DC.
        \pendente{resultado do controle de passo equivalente
        (\texttt{h2\_passo\_equivalente.yaml}).}
  \item \textbf{A fronteira de estabilidade aparece numa rede real?} Mapeamos a
        estabilidade no MNIST variando a taxa de aprendizado e o momento, e
        variando $K_d$ (inclusive negativo). Comparamos com a fronteira prevista pela
        curvatura máxima da função de erro na inicialização ($\lambda_{\max}$, o
        maior autovalor da matriz Hessiana).
        \begin{center}
        \figura{varredura_lr_momentum}{0.95}{Rode \texttt{sweep.py} com
        \texttt{mnist\_momentum\_lr.yaml} e depois \texttt{report\_assets.py -{}-sweep-lr}.}
        \end{center}
        \begin{center}
        \figura{varredura_kd}{0.95}{Rode \texttt{sweep.py} com
        \texttt{mnist\_pid\_kd.yaml} e depois \texttt{report\_assets.py -{}-sweep-kd}.}
        \end{center}
  \item \textbf{Testes estatísticos.} Teste $t$ pareado (mesmas sementes) e de Welch
        para as diferenças da Tab.~\ref{tab:fig4}, em especial PID com e sem o termo D.
\end{itemize}

% =============================================================================
\section{Como conferir os resultados}

Todos os números e figuras deste relatório podem ser verificados:
\begin{enumerate}
  \item \textbf{Origem.} Cada arquivo de resultado registra o \emph{commit} exato do
        código, se havia alterações não salvas, o hardware e a configuração. As
        tabelas geradas trazem essa origem num comentário no próprio \texttt{.tex}.
  \item \textbf{Auditoria.} O script \texttt{audit\_results.py} recalcula médias e
        desvios a partir das curvas de cada semente, confere se os valores são
        válidos e, opcionalmente, \emph{retreina} uma semente e compara época a época
        com o valor salvo.
  \item \textbf{Controle negativo.} Com taxa de aprendizado zero, os pesos não mudam,
        e a acurácia precisa ficar no nível do acaso (cerca de 10\,\%). Isso mostra que a
        acurácia reportada vem do treino, e não de um erro no cálculo
        (\texttt{controle\_negativo\_lr0.yaml}).
  \item \textbf{Comparação externa.} Nossos valores no MNIST são compatíveis com os
        registros publicados pelos autores (Tab.~\ref{tab:fig4}).
\end{enumerate}

% =============================================================================
\section{Próximos passos}

\begin{itemize}
  \item fechar as análises em andamento e os testes estatísticos;
  \item estender a mesma análise ao Adam e ao Adam-PID de Dai et al.\ (2023). Como a
        normalização do Adam é não linear, será preciso linearizar em torno de um
        ponto de operação;
  \item mapear sistematicamente as regiões de estabilidade nos hiperparâmetros
        (objetivo c da proposta) e relacioná-las com a capacidade de generalização.
\end{itemize}

% =============================================================================
\begin{thebibliography}{9}
\bibitem{an2018} W.~An, H.~Wang, Q.~Sun, J.~Xu, Q.~Dai e L.~Zhang.
  A PID controller approach for stochastic optimization of deep networks.
  In: \emph{IEEE/CVF Conference on Computer Vision and Pattern Recognition (CVPR)},
  p.~8522--8531, 2018. Código: \url{https://github.com/tensorboy/PIDOptimizer}.
\bibitem{dai2023} M.~Dai et al.
  PID controller-based adaptive gradient optimizer for deep neural networks.
  \emph{IET Control Theory \& Applications}, v.~17, n.~15, p.~2032--2037, 2023.
\bibitem{astrom} K.~J.~Åström e B.~Wittenmark.
  \emph{Adaptive Control}. 2.\ ed. Mineola: Dover, 2008.
\bibitem{wilson2017} A.~C.~Wilson et al.
  The marginal value of adaptive gradient methods in machine learning.
  \emph{Advances in Neural Information Processing Systems}, v.~30, 2017.
\end{thebibliography}

\end{document}
